% يُترجم باستخدام: xelatex أو lualatex (مرتين للمراجع إن وُجدت)
% مثال: xelatex "مسابقة-الجامعة-موعد.tex"
\documentclass[12pt,a4paper]{article}

\usepackage{fontspec}
\usepackage{geometry}
\geometry{margin=2.2cm}
\usepackage{polyglossia}
\setdefaultlanguage{arabic}
\setotherlanguage{english}

% خط عربي — غيّر الاسم إذا لم يُوجد على جهازك (مثلاً: Scheherazade New، Traditional Arabic، Noto Naskh Arabic)
\newfontfamily\arabicfont[Script=Arabic,Scale=1.05]{Amiri}
\renewcommand{\familydefault}{\arabicfont}

\usepackage{titlesec}
\usepackage{enumitem}
\setlist{nosep,leftmargin=*}
\usepackage{hyperref}
\hypersetup{colorlinks=true,linkcolor=blue!60!black,urlcolor=blue!60!black}

\title{\textbf{مشروع «موعد+»\\منصة رقمية لإدارة المواعيد والطوابير الطبية\\مع مساعد طبي ذكي لتوجيه المريض}}
\author{اسم الطالب / الفريق: \underline{\hspace{6cm}} \quad الرقم الجامعي: \underline{\hspace{3cm}}}
\date{\today}

\begin{document}
\maketitle
\vspace{-0.5em}
\begin{center}\small
\textbf{مستند تقديمي لمنافسة جامعية — وصف الفكرة والمشكلة والحل والميزة المقترحة}
\end{center}
\vspace{0.6em}

\section*{أولاً: المشكلة والحاجة}
تتزايد الضغوط على المنظومة الصحية مع ازدحام العيادات، وطول أوقات الانتظار، وصعوبة \textbf{اختيار التخصص المناسب} عند ظهور أعراض غير واضحة للمريض العادي، ما يؤدي أحياناً إلى تأخير المراجعة أو حجز موعد في تخصص غير ملائم. في الوقت نفسه، تبقى أنظمة الحجز التقليدية أو الجزئية الرقمية غالباً \textbf{لا تربط} بين فهم بسيط للأعراض وبين \textbf{أفضل الخيارات المتاحة فعلياً} من الأطباء المسجّلين في النظام حسب الجودة والتقييم.

\section*{ثانياً: الفكرة العامة للمشروع «موعد+»}
«موعد+» منصة تهدف إلى \textbf{تسهيل رحلة المريض} من البحث عن طبيب مناسب حتى \textbf{حجز موعد} وإدارة \textbf{طابور الانتظار} داخل العيادة، مع لوحة تحكم للطبيب لمتابعة المواعيد والطابور. يعتمد المشروع على بنية تقنية حديثة (تطبيق للمريض، وويب للطبيب) وقاعدة بيانات مركزية لحفظ بيانات الأطباء والمواعيد والتقييمات، بحيث تكون المعلومات \textbf{موحّدة وقابلة للاستعلام} لعرض الأطباء الأكثر تميزاً وفق معايير واضحة.

\section*{ثالثاً: مكوّنات النظام (باختصار)}
\begin{itemize}
  \item \textbf{تطبيق المريض:} تسجيل وتسجيل دخول، استعراض الأطباء والتخصصات، عرض التفاصيل والموقع، وحجز موعد ضمن أوقات متاحة، ومتابعة الطابور عند الحاجة.
  \item \textbf{لوحة الطبيب (ويب):} إدارة الملف والجدول، وقائمة المواعيد، وطابور اليوم، وإشعارات فورية عند التحديثات.
  \item \textbf{قاعدة البيانات والخدمات:} تخزين بيانات الأطباء (منها التخصص والتقييم وعدد المراجعات)، والمواعيد، وإعدادات الطابور، مع سياسات وصول تحكم من يقرأ ويعدّل ماذا.
\end{itemize}

\section*{رابعاً: الميزة المقترحة — «مساعد طبي ذكي» داخل التطبيق}
تُضاف طبقة تفاعلية جديدة تُكمّل الفكرة الأساسية ولا تحل محل الطبيب:

\begin{enumerate}
  \item \textbf{واجهة محادثة (شات)} داخل تطبيق المريض لاستيعاب وصف الأعراض بلغة يومية واضحة، مع توجيه أسئلة بسيطة عند الحاجة لزيادة الدقة دون إرهاق المستخدم.
  \item \textbf{تحليل بواسطة ذكاء اصطناعي} يُنتج \textbf{توجيهاً نحو تخصص طبي مناسب} كخطوة أولى، مع الالتزام بحد أمان صريح: \textbf{عدم تقديم تشخيص نهائي للمرض}، بل اقتراح «أي تخصص قد يكون الأنسب للمراجعة» وفق ما وصفه المستخدم.
  \item \textbf{ربط تلقائي بقاعدة البيانات:} بعد تحديد التخصص (أو أقرب تخصص مدعوم في النظام)، يُستدعى من الجدول نفسه قائمة بالأطباء في هذا التخصص، \textbf{مرتبين بحسب أعلى التقييمات} (مع معايير ثانوية مثل عدد المراجعات عند التعادل)، لضمان عرض \textbf{أفضل الخيارات المتاحة فعلياً} وليس اقتراحات عامة من الإنترنت.
  \item \textbf{قائمة مختصرة وقابلة للحجز:} يعرض المساعد للمريض عدداً محدوداً من الأطباء المقترحين مع \textbf{لمحة عن المواعيد أو الأوقات المتاحة} ضمن منطق الجدولة الحالي في التطبيق، لتسهيل \textbf{الحجز الفوري} أو الانتقال السريع إلى شاشة الحجز لنفس الطبيب.
\end{enumerate}

\section*{خامساً: الأمان والمسؤولية والقيمة المضافة}
يُفترض أن يُرفق المساعد بـ\textbf{تنويهات ثابتة}: أنه أداة مساعدة وليس بديلاً عن الفحص الطبي، وأن حالات الطوارئ تستوجب الاتصال بالجهات المختصة فوراً. القيمة المضافة للجامعة وللمسابقة تكمن في \textbf{دمج خدمة صحية رقمية عملية} مع \textbf{استخدام مسؤول للذكاء الاصطناعي} كجسر بين وصف الأعراض وتخصص حقيقي في النظام، ثم \textbf{ترشيح أطباء مُقيَّمين} وربط ذلك مباشرة بمسار الحجز—ما يجعل الفكرة \textbf{متكاملة من المشكلة إلى الحل التشغيلي}.

\vfill
\begin{center}\small
\textbf{ملاحظة للجنة:} يمكن إرفاق لقطات شاشة للتطبيق ولوحة الطبيب ولقطة مخطط تدفق (من المريض $\rightarrow$ الشات $\rightarrow$ التخصص $\rightarrow$ قائمة الأطباء $\rightarrow$ الحجز) عند العرض الشفهي أو الملحق البصري.
\end{center}

\end{document}
